\chapter{Protocole d'échange pour un cas métier concret}
\label{chap:protocole}

Pour illustrer la « pensée protocolaire » du projet, ce chapitre formalise l'échange d'un
cas métier complet --- la création d'un \texttt{Business} puis l'ajout d'un service à son
catalogue --- comme on spécifierait un protocole réseau : en-têtes obligatoires, payload,
réponses, codes d'erreur et sémantique de transport.

\section{Cas choisi : création d'un Business et ajout au catalogue}

Ce cas correspond à deux requêtes distinctes mais liées, typiques de l'enregistrement d'une
nouvelle activité (auto-école, cabinet de conseil, etc.) dans un tenant.

\section{Protocole : création d'un Business}

\noindent\textbf{Endpoint :} \texttt{POST /api/businesses}\\
\textbf{En-têtes obligatoires :}
\begin{itemize}
  \item \texttt{Authorization: Bearer <JWT>} --- authentification (analogie : ticket de
        session) ; le claim \texttt{tenantId} joue le rôle de l'adresse réseau.
\end{itemize}

\begin{lstlisting}[language=json,caption={Corps de la requête (payload)},captionpos=b]
 
{
  "name": "Auto-Ecole Centrale",
  "domain": "driving_school",
  "description": "Numero 1 dans la formation auto",
  "typeOfInvolvedActors": ["instructor", "student"],
  "requiredJobProfile": "moniteur",
  "tenantId": "3fa85f64-5717-4562-b3fc-2c963f66afa6"
}
\end{lstlisting}

\begin{lstlisting}[language=json,caption={Réponse (succès 201 Created)},captionpos=b]
{
  "id": "550e8400-e29b-41d4-a716-446655440000",
  "name": "Auto-Ecole Centrale",
  "domain": "driving_school",
  "createdAt": "2025-05-04T10:00:00Z"
}
\end{lstlisting}

\paragraph{Codes d'erreur (analogie : messages ICMP).}
\begin{itemize}
  \item \texttt{400 Bad Request} : payload invalide ;
  \item \texttt{401 Unauthorized} : JWT manquant ou invalide ;
  \item \texttt{403 Forbidden} : rôle insuffisant ;
  \item \texttt{404 Not Found} : ressource inexistante ;
  \item \texttt{409 Conflict} : ressource déjà existante.
\end{itemize}

\section{Protocole : ajout au catalogue de services}

\noindent\textbf{Endpoint :} \texttt{POST /api/service-catalogues/business/\{businessId\}}\\
\textbf{En-têtes :} identiques, plus une corrélation optionnelle.

\begin{lstlisting}[language=json,caption={Corps de la requête},captionpos=b]
{
  "name": "Cours de conduite",
  "description": "Formation theorique et pratique de 3 mois",
  "basePrice": 150000,
  "currency": "XAF",
  "isAvailable": true
}
\end{lstlisting}

\paragraph{Sémantique du transport.}
\begin{itemize}
  \item \textbf{Mode :} synchrone (REST) --- analogie TCP (fiable, ordonné) ;
  \item \textbf{Timeout :} de l'ordre de 30 secondes ;
  \item \textbf{Idempotence :} non (chaque \texttt{POST} crée une nouvelle ressource) ;
  \item \textbf{Retry :} recommandé avec \emph{backoff} exponentiel (au plus 3 tentatives).
\end{itemize}

% \section{Diagramme de séquence}

% \figplaceholder{Diagramme de séquence UML du cas « création Business + ajout au catalogue » :
% acteurs Provider, API Gateway, TenantFilter, BusinessService, ServiceCatalogueService, base
% de données ; messages REST et réponses 201, propagation du traceId.}{6.5cm}
% \captionof{figure}{Diagramme de séquence du cas métier « Business + Catalogue »}
% \label{fig:sequence}
